\documentclass[11pt]{article}
\usepackage{../homework}

\profname{Dr. Stanislaw Szarek}
\classcode{MATH 423 Measure Theory}
\assignnum{Assignment 12}
\duedate{23 November 2026}

\begin{document}
All problems from Folland, \textit{Real Analysis} (2e).

\begin{problem}{2.7.63}
  The technique used to prove Proposition 2.54 can also be used to integrate any polynomial over \( S^{n-1} \).
  In fact, suppose \( f(x) = \prod\limits_1^n x_j^{\alpha_j} \ (\alpha_j \in \mathbb{N} \cup \{0\}) \) is a monomial.
  Then \( \int f \dd{\sigma} = 0 \) if any \(\alpha_j\) is odd, and if all \(\alpha_j\)'s are even,
  \begin{equation}
  \label{eq:1}
  \int f \dd{\sigma} = \frac{2\Gamma(\beta_{1}) \cdots \Gamma(\beta_n)}{\Gamma(\beta_1 + \dots + \beta_{n})}, \text{ where } \beta_j = \frac{\alpha_{j}+1}{2}.
  \end{equation}
\end{problem}
\begin{proof}
  We consider the auxiliary integral \( \int f(x) e^{-|x|^2} \dd{x} \).
  First of all, defining \( g = f e^{-|x|^2}\), we have, after expanding the definition of \( f \), using \( |x|^2 = \sum\limits_1^n x_j^2 \), and appealing to Tonelli's theorem,
  \begin{equation}
    \label{eq:7}
    \begin{aligned}
      \int_{\mathbb{R}^{n}} |g| = \int \prod_1^n |x_{j}|^{\alpha_{j}} e^{-x_{j}^2} \dd{m} = \prod_1^n \int_{-\infty}^{\infty} |x_j|^{\alpha_j} e^{-x_j^2} \dd{x_{j}}.
    \end{aligned}
  \end{equation}
  Each factor is finite, since, writing \( t = x_j^2 \) and appealing to the symmetry,
  \begin{equation}
  \label{eq:8}
  \begin{aligned}
    \int |x_j|^{\alpha_j} e^{-x_j^2} \dd{x_j} &= 2 \int_0^{\infty} |x_j|^{\alpha_j} e^{-x_j^2} \dd{x_j} \\
    &= \int_0^{\infty} t^{(\alpha_j - 1)/2} e^{-t} \dd{t} = \Gamma\left(\frac{\alpha_{j} + 1}{2}\right) < \infty.
  \end{aligned}
\end{equation}
This shows that \( g \in L_1 \), so that for the remainder of the problem we may appeal to Fubini's theorem freely (note that in particular Theorem 2.49 will give integrability of \( f \) against the surface area measure.)

We can immediately eliminate the case where any of \(\alpha_j\) are odd.
Since the Lebesgue integral is invariant under reflection (\( u = -x \) gives \( \int h(x) = \int h(-x)  \) since the determinant of such a transformation has norm \( 1 \)), and since we can split the integral into a product of each component as in \cref{eq:7}, if any of the \(\alpha_j\) are odd, that component is an odd function.
Invariance under reflection then says that
\begin{equation}
\label{eq:9}
\int x_j^{\alpha_j} \dd{x_j} = \int (-x_j)^{\alpha_j},
\end{equation}
but since \( \alpha_j \) is odd, \( (-x_j)^{\alpha_j} = -(x_j)^{\alpha_j} \), so that both sides of \cref{eq:9} must vanish.
Therefore, we can consider only the case when all \( \alpha_j \) are even.

Now, by Theorem 2.49 we can also use a polar coordinate substitution so that
\begin{equation}
\label{eq:5}
\begin{aligned}
  \int f e^{-|x|^2} \dd{x} &= \int\limits_0^{\infty} \int\limits_{S^{n-1}} f(r x^{\prime}) r^{n-1} e^{-r^2} \dd{\sigma} \dd{r} \\
                        &= \int\limits_0^{\infty} \int\limits_{S^{n-1}} \prod_1^{n}(r x_j^{\prime})^{\alpha_j} r^{n-1} e^{-r^2} \dd{\sigma} \dd{r} \\
                        &= \int\limits_0^{\infty} \int\limits_{S^{n-1}} \left( \prod_1^n x^{\prime \alpha_j}_j \right) r^{\alpha_1 + \dots + \alpha_n + n - 1} e^{-r^2} \dd{\sigma} \dd{r} \\
                        &= \frac{1}{2} \int\limits_0^{\infty} t^{(\alpha_1 + \dots + \alpha_n + n)/2 - 1} e^{-t} \dd{t} \int f(x^{\prime}) \dd{\sigma} \\
                        &= \frac{1}{2} \Gamma((\alpha_1 + \dots + \alpha_n + n)/2) \int f \dd{\sigma} \\
                        &= \frac{1}{2} \Gamma(\beta_1 + \dots + \beta_n) \int f \dd{\sigma},
\end{aligned}
\end{equation}
where \( \beta_j \coloneq (\alpha_j + 1)/2 \).

Finally, returning to \cref{eq:8} and obeying the product gives (since for even \(\alpha_j\), the absolute value is superfluous)
\begin{equation}
\label{eq:10}
\int_{\mathbb{R}^2} f \dd{x} = \prod_1^n \Gamma(\beta_j),
\end{equation}

and therefore, dividing \cref{eq:5} leaves us with
\begin{equation}
\label{eq:11}
\int f \dd{\sigma} = \frac{2\Gamma(\beta_{1})\Gamma(\beta_2) \cdots \Gamma(\beta_n)}{\Gamma(\beta_1 + \Gamma(\beta_2) + \dots + \beta_{n})}.
\end{equation}
\end{proof}

\begin{prop}
  Let \(\nu\) be a signed measure on \( (X, \mathcal{M}) \).
  If \( \{E_j\} \) is an increasing sequence in \( \mathcal{M} \), then \( \nu(\bigcup_1^{\infty} E_j) = \lim_{j\rightarrow \infty} \nu(E_j)\).
  If \( \{E_j\} \) is a decreasing sequence in \( \mathcal{M} \) and \( \nu(E_1) \) is finite, then \( \nu(\bigcap_1^{\infty}) = \lim_{j\rightarrow\infty} \nu(E_j) \).
\end{prop}

\begin{problem}{3.1.1}
Prove Proposition 3.1.
\end{problem}
\begin{proof}
  We begin with the increasing case (continuity from below).
  Note that if, for some \( k \in \mathbb{N} \), \( \nu(E_k) = \infty\), then since the sequence increases we have \( \lim_{j\rightarrow \infty} \nu(E_j) = \infty = \nu(\bigcup_1^{\infty}E_j)\).
  It suffices to just consider the finite case.
  Just as in the positive measure case, set \( E_0 = \emptyset \).
  Then,
  \begin{equation}
  \label{eq:2}
  \begin{aligned}
    \nu(\bigcup_1^{\infty} E_j) &= \nu(\bigcup_1^{\infty} (E_j \setminus E_{j-1})) = \sum\limits_1^{\infty} \nu(E_j \setminus E_{j-1}) \\
    &= \lim_{N \rightarrow \infty} \sum\limits_1^N \nu(E_j \setminus E_{j-1}),
  \end{aligned}
\end{equation}
and since this sum is telescoping, as \( \nu(E_j \setminus E_{j-1}) = \nu(E_j) - \nu(E_{j-1}) \), we see that each of the partial sums is \( \sum\limits_1^n \nu(E_j \setminus E_{j-1}) = \nu(E_j) \), so that with \cref{eq:2} we get
\begin{equation}
\label{eq:3}
\nu(\bigcup_1^{\infty} E_j) = \lim_{j\rightarrow\infty} \nu(E_j).
\end{equation}

For the decreasing case (continuity from above), assuming \( \nu(E_1) \) is finite, the sequence \( F_j = E_1 \setminus E_j \) is an increasing sequence with \(\nu(F_j) = \nu(E_1) - \nu(E_j)\) and \( \bigcup_1^{\infty} F_j = E_1 \setminus \bigcap_1^{\infty} E_j\).
If any of \( E_{j} \) satisfy \( \nu(E_j) = -\infty \), then clearly the corresponding \( F_j \) will have \( \nu(F_j) = \infty \), in which case, just as for the first half, the remaining tail will all have infinite measure, so the result holds vacuously.

We can then appeal to the previous result to get that
\begin{equation}
\label{eq:12}
\begin{aligned}
  \nu(E_1) &= \nu(\bigcup_1^{\infty} F_j) + \nu(\bigcap_1^{\infty} E_j) = \lim_{j \rightarrow \infty} \nu(F_j) + \nu(\bigcap_1^{\infty} E_j) \\
  &= \lim_{j \rightarrow \infty} (\nu(E_1) - \nu(E_j)) + \nu(\bigcap_1^{\infty} E_j),
\end{aligned}
\end{equation}
and since by assumption \(\nu(E_1) \) is finite, subtracting that constant value from both sides gives
\begin{equation}
\label{eq:13}
\nu(\bigcap_1^{\infty} E_j) = \lim_{j\rightarrow \infty} \nu(E_j),
\end{equation}
as desired.
\end{proof}
\end{document}
